\documentclass[11pt]{article}

% ---------- packages ----------
\usepackage[utf8]{inputenc}
\usepackage[T1]{fontenc}
\usepackage{lmodern}
\usepackage[a4paper,margin=2.4cm]{geometry}
\usepackage{xcolor}
\usepackage{listings}
\usepackage{enumitem}
\usepackage{titlesec}
\usepackage{fancyhdr}
\usepackage{parskip}
\usepackage{booktabs}
\usepackage[hidelinks]{hyperref}

% ---------- palette (gonor.me) ----------
\definecolor{navy}{HTML}{1B2A4A}
\definecolor{terracotta}{HTML}{C17654}
\definecolor{amber}{HTML}{D4A574}
\definecolor{sage}{HTML}{7A8B6F}
\definecolor{cream}{HTML}{EDE6DD}
\definecolor{codebg}{HTML}{F5F0EA}
\definecolor{rule}{HTML}{D8CFC2}

\hypersetup{
  colorlinks=true,
  linkcolor=navy,
  urlcolor=terracotta,
  citecolor=terracotta
}

% ---------- headings ----------
\titleformat{\section}{\Large\bfseries\color{navy}}{\thesection}{0.6em}{}
\titleformat{\subsection}{\large\bfseries\color{terracotta}}{\thesubsection}{0.6em}{}
\titlespacing*{\section}{0pt}{1.4em}{0.5em}

\pagestyle{fancy}
\fancyhf{}
\lhead{\small\color{navy}SEO / AEO --- gonor.me}
\rhead{\small\color{navy}\thepage}
\renewcommand{\headrulewidth}{0.4pt}
\renewcommand{\headrule}{\hbox to\headwidth{\color{rule}\leaders\hrule height \headrulewidth\hfill}}

% ---------- code listings ----------
\lstset{
  basicstyle=\ttfamily\small,
  backgroundcolor=\color{codebg},
  frame=leftline,
  framerule=2pt,
  rulecolor=\color{amber},
  xleftmargin=1.2em,
  framexleftmargin=1.2em,
  breaklines=true,
  columns=fullflexible,
  keepspaces=true,
  showstringspaces=false,
  aboveskip=1em,
  belowskip=1em
}

% ---------- helper ----------
\newcommand{\term}[1]{\textcolor{navy}{\textbf{#1}}}

\begin{document}

% ================= TITLE =================
\begin{titlepage}
\centering
\vspace*{3cm}
{\color{navy}\Huge\bfseries SEO and AEO Optimization}\\[0.4cm]
{\color{terracotta}\Large A working guide from the \texttt{gonor.me} session}\\[1.2cm]
{\color{sage}\rule{0.5\textwidth}{1.5pt}}\\[1.2cm]
{\large What we changed, why it matters, and how to keep learning}\\[3cm]
{\large\textbf{Andr\'es Gonz\'alez Ortega}}\\[0.2cm]
{\normalsize Portfolio: \url{https://gonor.me}}\\[0.2cm]
{\normalsize July 2026}
\vfill
\end{titlepage}

\tableofcontents
\newpage

% ================= 1. CONTEXT =================
\section{Context}

This document explains the search-visibility work done on the portfolio site
\texttt{gonor.me} (an Astro static site) during one working session, and it
doubles as a self-study guide so the ideas stick.

The site was technically healthy (fast, no console errors, clean semantic
structure, mobile-safe). The problem was \emph{discoverability}: the entire
indexing layer pointed at the wrong domain, and the basic discovery files
(sitemap, robots, social image) were missing. In plain terms, the site worked
but search engines and answer engines were receiving contradictory signals
about what and where it was.

\subsection{Two goals, two audiences}

\begin{itemize}[leftmargin=1.4em]
  \item \term{SEO} (Search Engine Optimization): help \emph{search engines}
  (Google, Bing) crawl, understand, and rank your pages so a human clicks a
  blue link.
  \item \term{AEO} (Answer Engine Optimization): help \emph{answer engines}
  (ChatGPT, Perplexity, Claude, Google AI Overviews) find, trust, and
  \emph{cite} your content directly inside a generated answer.
\end{itemize}

They overlap heavily. Both need crawlable pages, a clear ``entity'' (who you
are), and machine-readable structure. AEO adds one demand: the text must be
\emph{extractable} --- a paragraph that answers a question and stands on its
own when quoted without the rest of the page.

% ================= 2. WHAT WE CHANGED =================
\section{What we changed this session}

The work split into four buckets: technical SEO, metadata and structured data,
AEO, and engagement. Each item below states \emph{the fix} and \emph{why it
matters}.

\subsection{Technical SEO (the plumbing)}

\begin{description}[leftmargin=0em, style=nextline]
  \item[Canonical domain unified to \texttt{gonor.me}.]
  Every page's \texttt{canonical}, \texttt{og:url}, JSON-LD \texttt{url}, and
  the sitemap host previously pointed at \texttt{gonorandres.github.io}, which
  in turn 301-redirects to \texttt{gonor.me}. That circular signal dilutes
  authority across two hosts and delays indexing. One config change fixed all
  of them at once.
  \item[Sitemap generated.] Added an i18n-aware \texttt{sitemap-index.xml}. A
  sitemap is the map you hand the crawler; without it, deep pages (28 projects,
  a blog, and artifacts in two languages) can take months to be discovered.
  \item[\texttt{robots.txt} created.] Tells crawlers what they may read and
  where the sitemap lives. It deliberately does \emph{not} block AI crawlers
  (GPTBot, ClaudeBot, PerplexityBot) --- being citable depends on them reading
  the site.
  \item[\texttt{hreflang} added.] The site has Spanish (\texttt{/}) and English
  (\texttt{/en/}) versions. \texttt{hreflang} tells Google they are
  translations of each other, not duplicate content, so the right language is
  shown to the right searcher.
  \item[404 marked \texttt{noindex}.] The error page should never appear in
  results; we told crawlers to skip it and dropped its canonical tag.
\end{description}

\begin{lstlisting}[caption={robots.txt --- allow everything, point at the sitemap}]
User-agent: *
Allow: /

Sitemap: https://gonor.me/sitemap-index.xml
\end{lstlisting}

\begin{lstlisting}[caption={hreflang pair emitted on every page}]
<link rel="alternate" hreflang="es" href="https://gonor.me/..." />
<link rel="alternate" hreflang="en" href="https://gonor.me/en/..." />
<link rel="alternate" hreflang="x-default" href="https://gonor.me/..." />
\end{lstlisting}

\subsection{Metadata and structured data}

\begin{description}[leftmargin=0em, style=nextline]
  \item[Titles and descriptions rewritten.] The \texttt{<title>} and meta
  description are the only text competing on the results page. We replaced a
  33-character name-only title with one that includes the terms people actually
  search (``Actuario y Cient\'ifico de Datos'', ``UNAM'') and wrote a full
  155-character description instead of a recycled 67-character one.
  \item[Person entity enriched.] The JSON-LD \texttt{Person} gained
  \texttt{sameAs} (GitHub + LinkedIn), \texttt{alumniOf} (UNAM),
  \texttt{knowsAbout}, \texttt{knowsLanguage}, and \texttt{email}. This is the
  backbone of AEO: it lets an engine merge the site, the GitHub, and the
  LinkedIn into one trusted entity.
  \item[Article schema completed.] Blog posts now emit \texttt{datePublished} /
  \texttt{dateModified} / \texttt{author} in both JSON-LD and Open Graph. Dates
  let Google show ``21 Mar 2026'' in the snippet and let answer engines judge
  whether the content is current.
  \item[Open Graph image created.] A \texttt{1200$\times$630} social card
  (\texttt{og-default.png}) was generated so links shared on LinkedIn,
  WhatsApp, or Slack render with an image instead of bare text.
\end{description}

\begin{lstlisting}[caption={The Person entity --- sameAs is the load-bearing field}]
{
  "@context": "https://schema.org",
  "@type": "Person",
  "name": "Andres Gonzalez Ortega",
  "url": "https://gonor.me",
  "jobTitle": "Actuario y Cientifico de Datos",
  "alumniOf": { "@type": "CollegeOrUniversity",
                "name": "Universidad Nacional Autonoma de Mexico" },
  "knowsAbout": ["Ciencia actuarial", "LISF / CUSF", "Machine Learning"],
  "sameAs": [ "https://github.com/GonorAndres",
              "https://www.linkedin.com/in/..." ]
}
\end{lstlisting}

\subsection{AEO --- being citable by answer engines}

\begin{description}[leftmargin=0em, style=nextline]
  \item[An ``About'' page with direct answers.] A new \texttt{/about/} page
  (both languages) is built from question headings (``Who is\dots?'', ``What
  does he do?'') where each answer's \emph{first sentence} is a complete,
  quotable statement. Answer engines lift self-contained sentences; scattered
  hero fragments are hard to cite.
  \item[\texttt{FAQPage} schema on flagship projects.] The two flagship posts
  got 3--4 real Q\&A pairs with \texttt{FAQPage} structured data --- exactly the
  questions a person asks a chatbot (``Where does the data come from?'').
  \item[\texttt{llms.txt} at the root.] An emerging convention: a short Markdown
  file summarizing who you are, what the site contains, and links to the key
  pages, written for language models.
\end{description}

\subsection{Engagement and trust}

\begin{itemize}[leftmargin=1.4em]
  \item Added \term{LinkedIn} to contact, footer, and the \texttt{sameAs} entity.
  \item \term{Self-hosted the CV} (\texttt{/docs/cv-\dots.pdf}) instead of a
  Google Drive link, so it never breaks and downloads can be measured as a GA4
  event.
  \item Wired \term{GA4 events} for the clicks that matter: \texttt{download\_cv},
  \texttt{click\_email}, \texttt{click\_repo}, \texttt{click\_demo}.
\end{itemize}

% ================= 3. GLOSSARY =================
\section{Key concepts, in one line each}

\begin{description}[leftmargin=0em, style=nextline]
  \item[Canonical URL] The one true address for a page. Prevents duplicate
  versions from splitting ranking signal.
  \item[\texttt{hreflang}] Declares language/region variants so the right
  translation is served and they are not treated as duplicates.
  \item[Sitemap] A machine list of every URL worth indexing; speeds discovery.
  \item[\texttt{robots.txt}] Crawl rules and the sitemap pointer. Allowing a bot
  is not the same as inviting it --- but blocking one guarantees invisibility.
  \item[Structured data / JSON-LD] A small script describing the page as data
  (\texttt{Person}, \texttt{Article}, \texttt{FAQPage}) using the
  \texttt{schema.org} vocabulary. This is how machines ``read'' meaning.
  \item[Entity \& \texttt{sameAs}] Your identity as a single node linked to all
  your profiles. The strongest AEO signal.
  \item[Open Graph (OG)] Tags that control the preview card when a link is
  shared on social platforms.
  \item[\texttt{FAQPage}] Structured Q\&A; eligible for rich results and easy
  for answer engines to quote.
  \item[\texttt{llms.txt}] A root Markdown file aimed at language models,
  summarizing the site.
  \item[Extractability] The property that a paragraph answers a question and
  survives being quoted alone. The core AEO writing habit.
\end{description}

% ================= 4. VERIFY =================
\section{How to verify the work}

Optimization you cannot measure is a guess. The tools:

\begin{itemize}[leftmargin=1.4em]
  \item \term{Google Search Console} --- register the domain, submit the
  sitemap, then watch what is indexed and which queries arrive. Without it you
  are blind. (\url{https://search.google.com/search-console})
  \item \term{Rich Results Test} --- validates your JSON-LD and shows which rich
  results a page is eligible for.
  (\url{https://search.google.com/test/rich-results})
  \item \term{Schema Markup Validator} --- \url{https://validator.schema.org}
  \item \term{Sharing debuggers} --- LinkedIn Post Inspector and Facebook
  Sharing Debugger re-scrape a URL to preview the OG card.
  \item \term{Local checks} --- \texttt{curl} for the raw HTML, or a headless
  browser to confirm \texttt{canonical}, \texttt{hreflang}, and JSON-LD render
  correctly. During this session every change was verified in the built HTML
  before shipping.
\end{itemize}

% ================= 5. LEARN MORE =================
\section{How to learn this properly}

A path from fundamentals to the frontier.

\subsection{Start with the primary sources}
\begin{itemize}[leftmargin=1.4em]
  \item \term{Google Search Central} --- the official documentation and SEO
  Starter Guide. The single best free resource.
  \url{https://developers.google.com/search/docs}
  \item \term{schema.org} --- the vocabulary reference for every structured-data
  type. Read the \texttt{Person}, \texttt{Article}, and \texttt{FAQPage} pages.
  \url{https://schema.org}
  \item \term{web.dev} (by Google) --- performance, Core Web Vitals, and
  accessibility, which feed ranking. \url{https://web.dev}
\end{itemize}

\subsection{Build the mental model}
\begin{itemize}[leftmargin=1.4em]
  \item Learn the crawl $\rightarrow$ index $\rightarrow$ rank pipeline. Most
  SEO problems are really \emph{crawl} or \emph{index} problems, not ranking
  ones.
  \item Learn the difference between \emph{on-page} (titles, headings, content,
  schema), \emph{technical} (speed, canonical, sitemap, rendering), and
  \emph{off-page} (links, mentions, entity reputation).
  \item Follow reputable, non-spammy sources: Google Search Central Blog,
  Ahrefs and Moz beginner guides, and Aleyda Solis's newsletter.
\end{itemize}

\subsection{Understand AEO, the new layer}
\begin{itemize}[leftmargin=1.4em]
  \item Read about \emph{entities} and the knowledge graph: search moved from
  matching keywords to understanding things. Your job is to be an
  unambiguous entity.
  \item Watch the \texttt{llms.txt} proposal and how AI crawlers
  (GPTBot, ClaudeBot, PerplexityBot) are documented; policy here changes fast.
  \item Practice the one writing habit that matters: open every page, section,
  and post with a self-contained answer. If a model quotes only your first two
  sentences, they should still make sense and be correct.
\end{itemize}

\subsection{Practice loop}
\begin{enumerate}[leftmargin=1.6em]
  \item Ship a change (a title, a schema block, a new answer paragraph).
  \item Validate it (Rich Results Test, Search Console URL inspection).
  \item Wait for data (impressions, clicks, indexed status).
  \item Read the query report; write the content those real queries want.
  \item Repeat. SEO/AEO is a compounding habit, not a one-time task.
\end{enumerate}

% ================= 6. CHECKLIST =================
\section{Quick reference checklist}

\begin{center}
\renewcommand{\arraystretch}{1.35}
\begin{tabular}{@{}p{0.62\textwidth} l@{}}
\toprule
\textbf{Item} & \textbf{Status this session} \\
\midrule
Canonical / og:url / JSON-LD on real domain & Done \\
\texttt{sitemap-index.xml} generated & Done \\
\texttt{robots.txt} (AI crawlers allowed) & Done \\
\texttt{hreflang} es / en / x-default & Done \\
Titles \& descriptions optimized per page & Done \\
\texttt{Person} entity with \texttt{sameAs} & Done \\
\texttt{Article} dates (published / modified) & Done \\
Open Graph image \texttt{1200$\times$630} & Done \\
\texttt{/about/} answer page (ES + EN) & Done \\
\texttt{FAQPage} on flagship projects & Done \\
\texttt{llms.txt} at root & Done \\
Self-hosted CV + GA4 click events & Done \\
\midrule
Register domain in Google Search Console & \emph{Your turn} \\
Submit the sitemap in Search Console & \emph{Your turn} \\
\texttt{www} subdomain DNS / redirect & \emph{Your turn} \\
\bottomrule
\end{tabular}
\end{center}

\vfill
\begin{center}
\small\color{sage}
Written as a companion to the \texttt{gonor.me} optimization session.\\
Verify claims against the primary sources above; search policy evolves.
\end{center}

\end{document}
